%%%%%%%%%%  scan idx 89  |  folio 76  %%%%%%%%%%
\SGAPrintPage{76}%
Formula (7) immediately gives \(y\lambda_{-1}(\check F)\in j^!(K(X'))\), provided that we have:
\sgaStableTarget{sga6.e00.eq.11}\begin{equation}\tag{11}
  y\,\gamma^{p-1}(E)\lambda^p(\check E)\in i^!(K(X)).
\end{equation}
Assume that this condition (11) is satisfied. By the corollary to lemma 2, \(\RR\) is true for \((Y',X',y\lambda_{-1}(\check F))\). On the other hand, to prove (2), it suffices, by formula (8), to prove that
\sgaStableTarget{sga6.e00.eq.12}\begin{equation}\tag{12}
  f^*\ch(i_!(y))
  \equiv
  f^*\bigl(i_*(\ch(y)T(E)^{-1})\bigr)
  \quad\text{mod.}\Ker f_*.
\end{equation}
Now, this last formula is verified immediately, using (6), (9), \(\RR\) for \((Y',X',y\lambda_{-1}(\check F))\) and the formula:
\sgaStableTarget{sga6.e00.eq.13}\begin{equation}\tag{13}
  \ch(\lambda_{-1}(\check F))=C^{p-1}(F)T(F)^{-1}\quad\text{(cf.\ formula (5)).}
\end{equation}

Thus, \(\RR\) is proved whenever \(y\,\gamma^{p-1}(E)=0\), in particular (since \(\gamma^{p-1}(E)\) has filtration \(\geq p-1\) \textsuperscript{*}in \(K(Y)\) filtered by the codimension of supports\textsubscript{*}) whenever \(p-1>\dim.Y\), i.e.\ whenever \(\dim.Y<\frac{\dim.X-1}{2}\). Now, we reduce to this case by combining lemma 1 and lemma 3 (taking for \(P\) a projective space of sufficiently large dimension \(r\), for example \(r>\dim.X+1\)). Q.E.D.

\sgaStableTarget{sga6.e00.rem.unnamed.occ6}N.B. We have used the fact that \(\gamma^{p-1}(E)\) has filtration \(\geq p-1\). We can prove this fact directly, without passing through Grassmannians as in \(\mathrm{R.R.R.}\) 2.16, by studying the bundle \(Y'\) of projective spaces associated with \(E\). This method gives, more generally, th.\ 2.14 of \(\mathrm{R.R.R.}\)

\sgaStableTarget{sga6.e00.extract-from-a-letter-of-grothendieck}\subsection*{Extract from a letter of Grothendieck}
\dots{} ``morally'', the new method rests on the determination of \(K(X)\) and \(A(X)\) when \(X\) is either a projective bundle associated with a vector bundle, or a variety obtained by blowing up a non%
%%%%%%%%%%  scan idx 90  |  folio 77  %%%%%%%%%%
\SGAPrintPage{77}%
singular subvariety. In this second case, I have not made the complete determination; a little something is missing. With the notation of the short paper, we must prove formula (9) (but as an \emph{equality}):
\sgaStableTarget{sga6.e00.eq.9.occ2}\begin{equation}\tag{9}
  f^*(i_*(\alpha))=j_*(\alpha C^{p-1}(F)),\qquad \alpha\in A(Y).
\end{equation}
We observe that the two sides have the same image under \(f_*\), and the same image under \(j^*\), so that the question is equivalent to:
\sgaStableTarget{sga6.e00.eq.9bis}\begin{equation}\tag{9 bis}
  \Ker f_*\cap\Ker j^*=0.
\end{equation}
In any case, (9) is true upon reduction to \(G(K(X'))\) (which destroys a little torsion, at most), since it then follows from (6); there is therefore hardly any doubt that it is true.

\dots I also point out to you that the \emph{equality} (9) follows very easily from (6) without passing through the troublesome (6 bis)\dots

\begin{flushleft}
p.c.c.J-P.S.
\end{flushleft}

\newpage

